% ============================================================
%  RAPPORT DE PFE — BrainScan AI
%  Auteure  : MOUMEN Doae Malak
%  Encadrante : Pr. Hasnâa CHAABI
%  Compilation : pdflatex rapport_pfe.tex  (×2 pour les refs)
% ============================================================

\documentclass[12pt, a4paper]{report}

% ── Encodage & langue ──────────────────────────────────────
\usepackage[utf8]{inputenc}
\usepackage[T1]{fontenc}
\usepackage[french]{babel}

% ── Mise en page ───────────────────────────────────────────
\usepackage[top=2.5cm, bottom=2.5cm, left=3cm, right=2.5cm]{geometry}
\usepackage{setspace}
\onehalfspacing

% ── Graphiques & couleurs ──────────────────────────────────
\usepackage{graphicx}
\usepackage[table,xcdraw]{xcolor}
\usepackage{tikz}
\usepackage{pgfplots}
\pgfplotsset{compat=1.18}

% ── Mathématiques ──────────────────────────────────────────
\usepackage{amsmath}
\usepackage{amssymb}
\usepackage{mathtools}

% ── Code source ────────────────────────────────────────────
\usepackage{listings}
\usepackage{inconsolata}   % police monospace lisible

\definecolor{codebg}{RGB}{245,247,250}
\definecolor{codecomment}{RGB}{106,153,85}
\definecolor{codestring}{RGB}{206,145,120}
\definecolor{codekeyword}{RGB}{86,156,214}
\definecolor{codenumber}{RGB}{181,181,181}

\lstdefinestyle{python}{
  language=Python,
  backgroundcolor=\color{codebg},
  basicstyle=\ttfamily\small,
  keywordstyle=\color{codekeyword}\bfseries,
  stringstyle=\color{codestring},
  commentstyle=\color{codecomment}\itshape,
  numberstyle=\tiny\color{codenumber},
  numbers=left,
  numbersep=8pt,
  stepnumber=1,
  showstringspaces=false,
  breaklines=true,
  frame=single,
  framerule=0.4pt,
  rulecolor=\color{gray!40},
  xleftmargin=15pt,
  framexleftmargin=15pt,
  tabsize=4,
  captionpos=b,
  literate=
    {é}{{\'e}}1 {è}{{\`e}}1 {ê}{{\^e}}1 {ë}{{\"e}}1
    {à}{{\`a}}1 {â}{{\^a}}1 {ù}{{\`u}}1 {û}{{\^u}}1
    {î}{{\^i}}1 {ô}{{\^o}}1 {ç}{{\c{c}}}1
    {É}{{\'E}}1 {È}{{\`E}}1 {Â}{{\^A}}1,
}
\lstset{style=python}

% ── Tableaux & listes ──────────────────────────────────────
\usepackage{booktabs}
\usepackage{tabularx}
\usepackage{enumitem}
\usepackage{multirow}
\usepackage{subcaption}   % sous-figures (a) / (b) avec légendes individuelles
\usepackage{float}        % placement [H] pour les figures non flottantes

% ── En-têtes & pieds de page ───────────────────────────────
\usepackage{fancyhdr}
\pagestyle{fancy}
\fancyhf{}
\fancyhead[L]{\small\leftmark}
\fancyhead[R]{\small BrainScan AI — PFE 2025}
\fancyfoot[C]{\thepage}
\renewcommand{\headrulewidth}{0.4pt}

% ── Liens hypertexte ───────────────────────────────────────
\usepackage[
  colorlinks = true,
  linkcolor  = blue!70!black,
  citecolor  = green!50!black,
  urlcolor   = cyan!60!black,
  pdftitle   = {Rapport PFE - BrainScan AI},
  pdfauthor  = {MOUMEN Doae Malak},
]{hyperref}

% ── Bibliographie ──────────────────────────────────────────
\usepackage[backend=biber, style=ieee, sorting=none]{biblatex}
\addbibresource{references.bib}   % créer ce fichier séparément

% ── Boîtes d'information ───────────────────────────────────
\usepackage{tcolorbox}
\tcbuselibrary{skins, breakable}

\newtcolorbox{infobox}[1][]{
  enhanced, breakable,
  colback=blue!5!white, colframe=blue!60!black,
  fonttitle=\bfseries, title=#1,
  left=6pt, right=6pt, top=4pt, bottom=4pt,
}
\newtcolorbox{warnbox}[1][]{
  enhanced, breakable,
  colback=orange!8!white, colframe=orange!70!black,
  fonttitle=\bfseries, title=#1,
  left=6pt, right=6pt, top=4pt, bottom=4pt,
}
\newtcolorbox{successbox}[1][]{
  enhanced, breakable,
  colback=green!6!white, colframe=green!55!black,
  fonttitle=\bfseries, title=#1,
  left=6pt, right=6pt, top=4pt, bottom=4pt,
}

% ── Commandes utilitaires ──────────────────────────────────
\newcommand{\code}[1]{\texttt{\small #1}}
\newcommand{\fichier}[1]{\texttt{\small\color{blue!70!black}#1}}

% ============================================================
%  DOCUMENT
% ============================================================
\begin{document}

% ── Numérotation romaine pour les pages préliminaires ──────
\pagenumbering{roman}
\setcounter{page}{0}

% ============================================================
%  PAGE DE GARDE
% ============================================================
\begin{titlepage}
  \begin{center}

    % Logo établissement (remplacez le chemin si nécessaire)
    % \includegraphics[width=3.5cm]{logo_ecole.png}\\[0.5cm]

    \rule{\linewidth}{1.5pt}\\[0.3cm]
    {\Large\scshape Université — École d'Ingénieurs}\\[0.1cm]
    {\large Département Informatique \& Intelligence Artificielle}\\
    \rule{\linewidth}{0.8pt}\\[1.5cm]

    {\huge\bfseries\color{blue!70!black}
      Détection Automatique de Tumeurs Cérébrales\\[0.4cm]
      par Deep Learning}\\[0.6cm]

    {\LARGE\itshape BrainScan AI}\\[0.3cm]
    {\large Application Web Django / ResNet50 / BraTS 2021}\\[2cm]

    \begin{minipage}{0.48\textwidth}
      \raggedright
      \textbf{Réalisé par :}\\[0.2cm]
      {\large MOUMEN Doae Malak}\\
      Spécialité : Data Science \& IA\\
      Promotion 2024--2025
    \end{minipage}
    \hfill
    \begin{minipage}{0.48\textwidth}
      \raggedleft
      \textbf{Encadrante :}\\[0.2cm]
      {\large Pr. Hasnâa CHAABI}\\
      Département IA \& Systèmes Intelligents
    \end{minipage}

    \vfill

    \rule{\linewidth}{0.8pt}\\[0.3cm]
    {\large Projet de Fin d'Études}\\[0.1cm]
    {\large Année universitaire 2024--2025}\\[0.3cm]
    \rule{\linewidth}{1.5pt}

  \end{center}
\end{titlepage}

% ============================================================
%  DÉDICACE
% ============================================================
\chapter*{Dédicace}
\addcontentsline{toc}{chapter}{Dédicace}
\thispagestyle{empty}

\begin{center}
  \itshape\large
  À ma famille, pour son soutien indéfectible tout au long de ce parcours.\\[0.5cm]
  À Pr. Hasnâa CHAABI, pour ses conseils éclairés et sa disponibilité.\\[0.5cm]
  À toute personne qui, un jour, bénéficiera de l'intelligence artificielle\\
  pour un diagnostic médical plus rapide et plus sûr.
\end{center}

% ============================================================
%  REMERCIEMENTS
% ============================================================
\chapter*{Remerciements}
\addcontentsline{toc}{chapter}{Remerciements}

Je tiens à exprimer ma profonde gratitude à \textbf{Pr. Hasnâa CHAABI},
encadrante de ce projet, dont les orientations scientifiques ont été
décisives pour la qualité du travail réalisé.

Mes remerciements vont également aux équipes de la compétition
\textit{BraTS 2021} (Brain Tumor Segmentation Challenge) pour la mise
à disposition du jeu de données d'IRM cérébrales annotées, et à la
communauté open-source TensorFlow/Keras pour les outils qui ont rendu
ce projet possible.

% ============================================================
%  RÉSUMÉ
% ============================================================
\chapter*{Résumé}
\addcontentsline{toc}{chapter}{Résumé}

Ce projet de fin d'études présente \textbf{BrainScan AI}, une application
web de détection automatique de tumeurs cérébrales à partir d'IRM en
format NIfTI (\code{.nii.gz}). Le système s'appuie sur une architecture
\textbf{ResNet50} pré-entraînée sur ImageNet, affinée (fine-tuning sur
les 20 dernières couches) avec une fonction de perte combinant
\textit{Focal Loss} et \textit{Label Smoothing} pour gérer
l'incertitude inhérente aux données médicales.

Le pipeline complet — du prétraitement par Z-Score dataset-level jusqu'à
la présentation du diagnostic à l'utilisateur — est intégré dans une
application \textbf{Django}. Une attention particulière a été portée à
l'expérience utilisateur : les probabilités brutes sont masquées au profit
d'un indicateur de \textit{fiabilité du diagnostic}, et un seuil de
décision optimal est calculé automatiquement via l'indice de Youden.

\textbf{Mots-clés :} Deep Learning, Détection de tumeurs, IRM, ResNet50,
Focal Loss, Django, BraTS 2021, Seuil de Youden.

\vspace{1cm}

\textbf{Abstract —} This final year project presents BrainScan AI, a web
application for automatic brain tumor detection from NIfTI MRI scans.
The system combines a fine-tuned ResNet50 backbone with a custom
Focal + Label Smoothing loss, deployed via a Django interface that
prioritises clinical reassurance over raw probability display.

% ============================================================
%  TABLE DES MATIÈRES
% ============================================================
\tableofcontents
\listoffigures
\listoftables
\lstlistoflistings

% ── Numérotation arabe dès le premier chapitre ─────────────
\cleardoublepage
\pagenumbering{arabic}
\setcounter{page}{1}

% ============================================================
%  CHAPITRE 1 — INTRODUCTION GÉNÉRALE
% ============================================================
\chapter{Introduction Générale}
\label{chap:intro}

\section{Contexte et Motivation}

Les tumeurs cérébrales représentent l'une des pathologies les plus
redoutables en oncologie. Selon l'Organisation Mondiale de la Santé,
les glioblastomes multiformes affichent un taux de survie à cinq ans
inférieur à 10~\%. La rapidité du diagnostic est un facteur pronostique
déterminant~: plus la lésion est détectée tôt, plus les options
thérapeutiques sont nombreuses.

L'Imagerie par Résonance Magnétique (IRM) constitue le gold standard
pour la détection et la caractérisation des tumeurs cérébrales.
Cependant, l'interprétation manuelle de centaines de coupes par examen
est une tâche longue, subjective et susceptible d'erreurs, surtout en
l'absence d'un radiologue spécialisé.

L'essor du \textit{deep learning}, et notamment des réseaux de neurones
convolutifs (CNN), offre une opportunité inédite~: automatiser la
détection des lésions à partir des volumes IRM, avec une précision
comparable à celle des experts humains.

\section{Objectifs du Projet}

Ce projet vise à concevoir et déployer un système de détection de
tumeurs cérébrales accessible via un navigateur web, avec les objectifs
suivants~:

\begin{enumerate}[leftmargin=*, label=\arabic*.]
  \item Entraîner un modèle CNN performant sur le dataset \textbf{BraTS 2021}.
  \item Intégrer des mécanismes anti-surconfiance (Focal Loss, Label Smoothing).
  \item Calculer automatiquement un seuil de décision cliniquement sûr.
  \item Développer une interface web rassurante pour le patient et le clinicien.
  \item Garantir la cohérence entre le prétraitement à l'entraînement et
        celui à l'inférence.
\end{enumerate}

\section{Organisation du Rapport}

Ce rapport est structuré en cinq chapitres~:
\begin{itemize}
  \item \textbf{Chapitre 2} — État de l'art en détection de tumeurs par DL.
  \item \textbf{Chapitre 3} — Méthodologie technique (modèle, loss, normalisation).
  \item \textbf{Chapitre 4} — Analyse, calibration et correction des anomalies.
  \item \textbf{Chapitre 5} — Réalisation de l'interface et expérience utilisateur.
  \item \textbf{Chapitre 6} — Conclusion et perspectives.
\end{itemize}

% ============================================================
%  CHAPITRE 2 — ÉTAT DE L'ART
% ============================================================
\chapter{État de l'Art}
\label{chap:etat_art}

\section{Détection de Tumeurs Cérébrales par Deep Learning}

Depuis les travaux fondateurs de Havaei et al. (2017) sur la
segmentation pixel-à-pixel, la détection de tumeurs cérébrales par
CNN a connu une évolution rapide. On distingue deux grandes approches~:

\begin{description}
  \item[Segmentation (U-Net, nnU-Net)] Prédit un masque binaire
        \(\hat{M}_{ij} \in \{0, 1\}\) pour chaque voxel.
        Très précise mais nécessite des annotations de segmentation
        complètes, coûteuses à produire.

  \item[Classification par coupe] Prédit un label binaire
        \(\hat{y} \in \{0, 1\}\) par coupe axiale.
        Plus simple, entraînable avec des labels de présence/absence,
        et suffisant pour un premier triage clinique.
\end{description}

Ce projet adopte la \textbf{classification par coupe} en raison de la
disponibilité des labels et de la simplicité de déploiement.

\section{Architectures de Référence}

\begin{table}[h]
  \centering
  \caption{Comparaison des architectures CNN pour la classification médicale}
  \label{tab:architectures}
  \begin{tabularx}{\textwidth}{lXrr}
    \toprule
    \textbf{Architecture} & \textbf{Caractéristiques} &
    \textbf{Param.} & \textbf{AUC typique} \\
    \midrule
    VGG-16     & Simple, profond, pas de skip-connections & 138M & 0.91 \\
    InceptionV3 & Multi-échelle, efficace en mémoire & 24M & 0.93 \\
    \textbf{ResNet50} & Skip-connections, fine-tuning robuste & 25M & \textbf{0.95} \\
    EfficientNetB4 & Scaling uniforme, SOTA récent & 19M & 0.96 \\
    \bottomrule
  \end{tabularx}
\end{table}

\textbf{ResNet50} a été retenu pour ce projet en raison de son excellent
compromis performance/taille, de sa disponibilité sur Keras avec les
poids ImageNet, et de sa robustesse éprouvée en imagerie médicale.

\section{Problématiques Spécifiques aux Données Médicales}

\subsection{Déséquilibre des Classes}

Dans le dataset BraTS 2021, de nombreuses coupes axiales d'un même
patient tumoral ne contiennent pas de pixels tumoraux (coupe au-dessus
ou en-dessous de la masse). Ce phénomène crée un déséquilibre naturel~:

\begin{equation}
  \frac{N_{\text{sain}}}{N_{\text{tumoral}}} \approx 3{,}5
\end{equation}

Un modèle non corrigé tend alors à classer systématiquement
\textit{sain} pour maximiser l'accuracy globale.

\subsection{Surconfiance des Réseaux de Neurones}

Les CNN modernes produisent des probabilités de sortie proche de 0 ou 1
(\textit{overconfidence}), rendant difficile l'interprétation clinique.
Cette propriété est particulièrement dangereuse en médecine~: un modèle
qui dit «~96~\% sain~» avec une vraie tumeur est plus dangereux qu'un
modèle qui dit «~70~\% sain~».

% ============================================================
%  CHAPITRE 3 — MÉTHODOLOGIE TECHNIQUE
% ============================================================
\chapter{Méthodologie Technique}
\label{chap:methodo}

\section{Vue d'Ensemble du Pipeline}

\begin{figure}[h]
  \centering
  \begin{tikzpicture}[
    node distance=0.7cm and 1.2cm,
    box/.style={
      draw, rounded corners=4pt,
      fill=blue!10, minimum height=0.9cm,
      minimum width=3.2cm, align=center,
      font=\small
    },
    arrow/.style={->, thick, blue!60!black}
  ]
    \node[box] (nii)    {Fichier .nii.gz};
    \node[box, right=of nii]  (slice)  {Coupe axiale\\(slice 75)};
    \node[box, right=of slice] (zscore) {Z-Score\\Dataset};
    \node[box, right=of zscore] (clahe)  {CLAHE};
    \node[box, below=of clahe] (resize) {Resize\\224×224};
    \node[box, left=of resize]  (rgb)   {Grayscale→RGB\\+ImageNet norm};
    \node[box, left=of rgb]     (model) {ResNet50\\(fine-tuned)};
    \node[box, left=of model]   (inv)   {Inversion\\1-p};
    \node[box, below=of inv]    (seuil) {Seuil Youden\\(auto)};
    \node[box, right=of seuil]  (ux)    {Interface\\Django};

    \draw[arrow] (nii)   -- (slice);
    \draw[arrow] (slice) -- (zscore);
    \draw[arrow] (zscore)-- (clahe);
    \draw[arrow] (clahe) -- (resize);
    \draw[arrow] (resize)-- (rgb);
    \draw[arrow] (rgb)   -- (model);
    \draw[arrow] (model) -- (inv);
    \draw[arrow] (inv)   -- (seuil);
    \draw[arrow] (seuil) -- (ux);
  \end{tikzpicture}
  \caption{Pipeline complet de BrainScan AI, de l'IRM brute au verdict affiché}
  \label{fig:pipeline}
\end{figure}

\section{Prétraitement — Normalisation Z-Score Dataset}
\label{sec:zscore}

\subsection{Problème : Hétérogénéité des Intensités IRM}

Les IRM cérébrales présentent des variations d'intensité importantes
d'un patient à l'autre (liées à la machine, au protocole d'acquisition,
à l'âge du patient). Une normalisation purement par image amplifie
artificiellement les images peu contrastées, ce qui nuit à la
généralisation du modèle.

\subsection{Solution : Z-Score au Niveau du Dataset}

La normalisation est calculée \textbf{une seule fois} sur l'ensemble
d'entraînement, puis réappliquée identiquement à la validation, au
test et à l'inférence en production.

\begin{equation}
  \label{eq:zscore}
  I'_{ij} = \frac{I_{ij} - \mu_{\text{train}}}{\sigma_{\text{train}} + \varepsilon}
\end{equation}

où \(\mu_{\text{train}}\) et \(\sigma_{\text{train}}\) sont calculés
sur les pixels cerveau (\(I_{ij} > 0\)) de toutes les coupes
d'entraînement.

Ces statistiques sont persistées dans le fichier
\fichier{dataset\_stats.json}~:

\begin{lstlisting}[caption={Contenu de dataset\_stats.json},
                   label={lst:stats_json}]
{
  "mean":      1729.70,
  "std":       16289.58,
  "n_images":  875,
  "threshold": 0.5312    # ajoute automatiquement apres entrainement
}
\end{lstlisting}

\subsection{Implémentation dans brats\_classifier.py}

\begin{lstlisting}[caption={Calcul des statistiques dataset (brats\_classifier.py)},
                   label={lst:compute_stats}]
def compute_dataset_stats(paths, slices, save_path=STATS_FILE):
    all_pixels = []
    for path, z in zip(paths, slices):
        s = extract_raw_slice(path, int(z))
        brain_px = s[s > 0]           # pixels cerveau uniquement
        if brain_px.size > 0:
            all_pixels.append(brain_px)

    all_pixels = np.concatenate(all_pixels)
    stats = {
        "mean": float(np.mean(all_pixels)),
        "std":  float(np.std(all_pixels)),
        "n_images": len(paths),
    }
    with open(save_path, "w") as f:
        json.dump(stats, f, indent=2)
    return stats
\end{lstlisting}

\subsection{Réapplication Cohérente en Production}

\begin{lstlisting}[caption={Z-Score dataset dans views.py (inférence)},
                   label={lst:zscore_infer}]
def _preprocess(slice_2d):
    stats = _get_stats()           # charge dataset_stats.json

    if stats is not None:
        mu, sigma  = stats["mean"], stats["std"]
        normalized = (slice_2d - mu) / (sigma + 1e-8)
    else:
        # fallback : Z-Score par image (moins robuste)
        nonzero = slice_2d[slice_2d > 0]
        mu      = nonzero.mean() if nonzero.size > 0 else 0.0
        sigma   = nonzero.std()  if nonzero.size > 0 else 1.0
        normalized = (slice_2d - mu) / (sigma + 1e-8)
\end{lstlisting}

\begin{infobox}[Invariant train/inférence]
  Le même fichier \fichier{dataset\_stats.json} est utilisé à
  l'entraînement et en production. Si ce fichier est absent ou
  corrompu, Django bascule automatiquement sur le Z-Score par image
  et journalise un avertissement.
\end{infobox}

\subsection{Vérification : CLAHE ne Lisse Pas les Tumeurs}

Un risque potentiel est que la normalisation globale suivie du CLAHE
(Contrast Limited Adaptive Histogram Equalization) atténue les zones
hypersignal caractéristiques des tumeurs.

\textbf{Analyse}~: le clip à \(\pm 3\sigma\) garantit que les valeurs
extrêmes (zones tumorales souvent au-dessus de \(+2\sigma\)) sont
préservées à leur maximum avant CLAHE. Le CLAHE, appliqué localement
avec \code{clipLimit=2.0}, augmente le contraste des structures subtiles
\textit{sans écraser} les zones déjà contrastées. Cette combinaison
\textbf{améliore} la visibilité des tumeurs pour le CNN.

\section{Fonction de Perte — Focal Loss + Label Smoothing}
\label{sec:loss}

\subsection{Motivations}

Deux problèmes distincts sont adressés par la fonction de perte
personnalisée~:

\begin{enumerate}
  \item \textbf{Déséquilibre des classes} : les exemples faciles
        (cerveau clairement sain) dominent le gradient et empêchent
        le modèle d'apprendre les cas ambigus.
  \item \textbf{Surconfiance} : la Binary Cross-Entropy standard pousse
        les probabilités vers 0 ou 1, produisant des prédictions
        extrêmes difficiles à interpréter cliniquement.
\end{enumerate}

\subsection{Formulation Mathématique}

\subsubsection{Binary Cross-Entropy avec Label Smoothing}

Soit \(y \in \{0, 1\}\) le label réel et \(\hat{y}\) la prédiction
sigmoïde. Avec un coefficient de lissage \(\varepsilon = 0.1\)~:

\begin{equation}
  \label{eq:smooth_label}
  \tilde{y} = y(1 - \varepsilon) + \frac{\varepsilon}{2}
  = \begin{cases}
      0.05 & \text{si } y = 0 \\
      0.95 & \text{si } y = 1
    \end{cases}
\end{equation}

La BCE avec label smoothing devient~:

\begin{equation}
  \label{eq:bce_smooth}
  \mathcal{L}_{\text{BCE}}^{\varepsilon}(\tilde{y}, \hat{y})
  = -\tilde{y}\log\hat{y} - (1-\tilde{y})\log(1-\hat{y})
\end{equation}

L'effet direct est de plafonner la pénalité maximale~: le modèle n'est
jamais récompensé pour avoir prédit exactement 0 ou 1, ce qui réduit
la surconfiance.

\subsubsection{Focal Loss}

Pour un exemple d'entraînement, la probabilité «~correcte~» est~:

\begin{equation}
  p_t = y\hat{y} + (1-y)(1-\hat{y})
\end{equation}

La Focal Loss introduit un terme de pondération exponentiel~:

\begin{equation}
  \label{eq:focal}
  w_{\text{focal}} = \alpha \cdot (1 - p_t)^{\gamma}
\end{equation}

\begin{itemize}
  \item \(\gamma = 2\)~: les exemples faciles (\(p_t \to 1\)) reçoivent
        un poids quasi nul ; les cas difficiles (\(p_t \to 0\)) reçoivent
        un poids plein.
  \item \(\alpha = 0.25\)~: facteur d'équilibre global.
\end{itemize}

\subsubsection{Combinaison Finale}

\begin{equation}
  \label{eq:combined_loss}
  \boxed{
    \mathcal{L}_{\text{total}} = \mathbb{E}\bigl[w_{\text{focal}}\bigr]
    \cdot \mathcal{L}_{\text{BCE}}^{\varepsilon}
  }
\end{equation}

\subsection{Implémentation}

\begin{lstlisting}[caption={Fonction de perte combinée (brats\_classifier.py)},
                   label={lst:loss}]
def build_loss(gamma=2.0, alpha=0.25, label_smoothing=0.1):
    bce_smooth = tf.keras.losses.BinaryCrossentropy(
        label_smoothing=label_smoothing,
        from_logits=False,
    )

    def focal_smooth_loss(y_true, y_pred):
        bce = bce_smooth(y_true, y_pred)

        y_pred_c = tf.clip_by_value(y_pred, 1e-7, 1 - 1e-7)
        p_t      = y_true * y_pred_c + (1 - y_true) * (1 - y_pred_c)
        focal_w  = alpha * tf.pow(1.0 - p_t, gamma)

        return tf.reduce_mean(focal_w) * bce

    return focal_smooth_loss
\end{lstlisting}

\begin{infobox}[Chargement en production avec custom\_objects]
  Keras ne sait pas retrouver \code{focal\_smooth\_loss} à partir du
  nom seul. Il faut la fournir explicitement lors du chargement~:
\begin{lstlisting}
focal_smooth_loss = _build_focal_smooth_loss(
    gamma=2.0, alpha=0.25, label_smoothing=0.1
)
model = load_model(
    model_path,
    custom_objects={"focal_smooth_loss": focal_smooth_loss},
)
\end{lstlisting}
\end{infobox}

\section{Architecture du Modèle}
\label{sec:archi}

\subsection{ResNet50 — Fine-Tuning Sélectif}

\begin{lstlisting}[caption={Construction du modèle ResNet50 (brats\_classifier.py)},
                   label={lst:model}]
def build_model(input_shape=(224, 224, 3)):
    inputs = layers.Input(shape=input_shape)

    # Augmentation integree (train uniquement, desactivee a l'inference)
    x = layers.RandomFlip("horizontal")(inputs)
    x = layers.RandomRotation(0.10)(x)

    base = ResNet50(weights="imagenet", include_top=False,
                    input_shape=input_shape)
    for layer in base.layers[:-20]:   # geler toutes sauf les 20 dernieres
        layer.trainable = False
    for layer in base.layers[-20:]:
        layer.trainable = True

    x = base(x, training=False)
    x = layers.GlobalAveragePooling2D()(x)

    x = layers.Dense(256, activation="relu",
                     kernel_regularizer=regularizers.l2(1e-3))(x)
    x = layers.BatchNormalization()(x)
    x = layers.Dropout(0.50)(x)

    x = layers.Dense(64, activation="relu",
                     kernel_regularizer=regularizers.l2(1e-3))(x)
    x = layers.BatchNormalization()(x)
    x = layers.Dropout(0.50)(x)

    output = layers.Dense(1, activation="sigmoid",
                          name="tumor_prob")(x)
    return Model(inputs=inputs, outputs=output)
\end{lstlisting}

\subsection{Poids des Classes Asymétriques}

Pour pénaliser davantage les faux négatifs (tumeurs manquées), des
poids asymétriques sont appliqués~:

\begin{equation}
  w_{\text{sain}} = 1.0 \qquad w_{\text{tumoral}} = 2.0
\end{equation}

\begin{lstlisting}[caption={Poids des classes dans model.fit()},
                   label={lst:classweights}]
class_weights = {0: 1.0, 1: 2.0}

history = model.fit(
    train_gen,
    validation_data=val_gen,
    epochs=50,
    callbacks=callbacks,
    class_weight=class_weights,
    verbose=1,
)
\end{lstlisting}

% ============================================================
%  CHAPITRE 4 — ANALYSE ET CALIBRATION
% ============================================================
\chapter{Analyse et Calibration du Modèle}
\label{chap:calibration}

\section{Correction de l'Inversion des Classes}
\label{sec:inversion}

\subsection{Symptôme Observé}

Lors des premiers tests en production, des IRM avec des masses tumorales
évidentes recevaient des probabilités proches de 0, tandis que des
coupes saines obtenaient des probabilités proches de 1. Le modèle
fonctionnait correctement, mais ses sorties étaient \textbf{sémantiquement
inversées}.

\subsection{Diagnostic}

L'inversion est due à l'ordre des classes lors de la collecte des
données sur Kaggle. La fonction \code{collect\_samples\_simple} a
traité les fichiers dans l'ordre alphabétique, plaçant les cerveaux
\textit{healthy} avant les tumeurs dans le générateur, résultant en~:

\begin{equation}
  \text{Classe~0} \leftarrow \text{Tumeur} \qquad
  \text{Classe~1} \leftarrow \text{Sain}
\end{equation}

au lieu de la convention attendue (0=Sain, 1=Tumeur).

\subsection{Correction Adoptée}

La correction la plus propre, qui ne nécessite pas de réentraîner
le modèle, est d'inverser la probabilité brute côté inférence~:

\begin{equation}
  \label{eq:inversion}
  \boxed{p_{\text{tumeur}} = 1 - p_{\text{brute}}}
\end{equation}

\begin{lstlisting}[caption={Inversion dans views.py (detector/views.py)},
                   label={lst:inversion}]
raw_output = model.predict(tensor, verbose=0)

# Le modele a ete entraine avec 0=Tumeur, 1=Sain
# --> inversion pour que tumor_prob = P(Tumeur)
if raw_output.shape[-1] == 1:
    tumor_prob = 1.0 - float(raw_output[0][0])   # sigmoidee
else:
    tumor_prob = 1.0 - float(raw_output[0][1])   # softmax classe 1
\end{lstlisting}

\begin{warnbox}[Bonne pratique pour les projets futurs]
  Pour éviter ce type d'inversion, il est recommandé de~:
  \begin{enumerate}
    \item Vérifier l'ordre des classes avec \code{np.unique(labels)}
          avant l'entraînement.
    \item Sauvegarder le mapping \code{\{"0": "sain", "1": "tumeur"\}}
          dans \fichier{dataset\_stats.json}.
    \item Valider sur un exemple connu avant tout déploiement.
  \end{enumerate}
\end{warnbox}

\section{Calcul du Seuil Optimal — Indice de Youden}
\label{sec:youden}

\subsection{Problématique du Seuil Fixe à 0.98}

Un seuil de décision fixe à \(0.98\) avait été initialement utilisé
pour éliminer les faux positifs. Cette approche s'est révélée trop
conservatrice~: des IRM présentant des masses tumorales claires avec
\(p_{\text{tumeur}} = 0.85\) étaient classées \textit{saines}.

\textbf{Le vrai problème} n'était pas le seuil, mais l'inversion des
classes (§\ref{sec:inversion}). Une fois l'inversion corrigée, le seuil
peut être abaissé à une valeur cliniquement raisonnée.

\subsection{Justification Mathématique — Indice de Youden}

L'indice de Youden \(J\) mesure la capacité globale du classifieur
à distinguer les deux classes~:

\begin{equation}
  \label{eq:youden}
  J(\tau) = \underbrace{\frac{TP}{TP + FN}}_{\text{Sensibilité}}
           + \underbrace{\frac{TN}{TN + FP}}_{\text{Spécificité}}
           - 1
\end{equation}

Le seuil optimal maximise cet indice~:

\begin{equation}
  \label{eq:youden_opt}
  \tau^* = \arg\max_{\tau \in [0,1]} J(\tau)
         = \arg\max_{\tau} \bigl[\text{TPR}(\tau) - \text{FPR}(\tau)\bigr]
\end{equation}

Ce seuil correspond au point de la courbe ROC le plus éloigné de la
diagonale aléatoire, tel qu'illustré à la
Figure~\ref{fig:roc_youden}.

\begin{figure}[h]
  \centering
  \begin{tikzpicture}
    \begin{axis}[
      width=8cm, height=7cm,
      xlabel={Taux de Faux Positifs (FPR)},
      ylabel={Taux de Vrais Positifs (TPR)},
      xmin=0, xmax=1, ymin=0, ymax=1,
      grid=both, grid style={line width=0.2pt, gray!30},
      legend pos=south east,
      title={Courbe ROC et point de Youden},
    ]
      % Diagonale aleatoire
      \addplot[dashed, gray, thick] coordinates {(0,0) (1,1)};
      \addlegendentry{Classifieur aléatoire}

      % Courbe ROC simulee
      \addplot[blue!70!black, thick, smooth] coordinates {
        (0,0) (0.05,0.45) (0.10,0.65) (0.15,0.75)
        (0.25,0.85) (0.40,0.92) (0.60,0.96) (0.80,0.98) (1,1)
      };
      \addlegendentry{Courbe ROC (AUC $\approx 0.93$)}

      % Point de Youden
      \addplot[red, only marks, mark=*, mark size=3pt]
        coordinates {(0.10, 0.65)};
      \addlegendentry{Seuil de Youden $\tau^*$}

      % Ligne verticale du point de Youden
      \addplot[red, dashed, thin] coordinates {(0.10,0) (0.10,0.65)};
      \addplot[red, dashed, thin] coordinates {(0,0.65) (0.10,0.65)};
    \end{axis}
  \end{tikzpicture}
  \caption{Point de Youden sur la courbe ROC — maximise sensibilité + spécificité}
  \label{fig:roc_youden}
\end{figure}

\subsection{Implémentation Automatique}

\begin{lstlisting}[caption={Calcul du seuil optimal (brats\_classifier.py)},
                   label={lst:youden}]
def find_optimal_threshold(model, val_gen,
                            method="youden",
                            stats_path=STATS_FILE):
    y_true, y_probs = [], []
    for Xb, yb in val_gen:
        y_probs.extend(model.predict(Xb, verbose=0).flatten())
        y_true.extend(yb)
    y_true  = np.array(y_true)
    y_probs = np.array(y_probs)

    fpr, tpr, thresholds = roc_curve(y_true, y_probs)
    youden   = tpr - fpr                  # eq. (4.4)
    best_idx = int(np.argmax(youden))
    best_t   = float(thresholds[best_idx])

    # Clamp cliniquement raisonnable
    best_t = float(np.clip(best_t, 0.25, 0.75))

    # Persistance dans dataset_stats.json
    with open(stats_path) as f:
        stats = json.load(f)
    stats["threshold"] = best_t
    with open(stats_path, "w") as f:
        json.dump(stats, f, indent=2)

    return best_t
\end{lstlisting}

\subsection{Lecture Dynamique dans Django}

\begin{lstlisting}[caption={Chargement du seuil optimal depuis JSON (views.py)},
                   label={lst:threshold_load}]
_raw_stats = _get_stats()
DECISION_THRESHOLD = (
    float(_raw_stats["threshold"])
    if _raw_stats and "threshold" in _raw_stats
    else 0.50
)
has_tumor = tumor_prob >= DECISION_THRESHOLD
\end{lstlisting}

\begin{successbox}[Avantages de cette approche]
  \begin{itemize}
    \item Le seuil s'adapte automatiquement à chaque nouveau modèle
          sans modifier le code Django.
    \item La valeur est calmée dans $[0.25, 0.75]$, ce qui garantit
          un comportement cliniquement interprétable.
    \item En cas de fichier absent, Django utilise le fallback $0.50$
          (règle d'or de la probabilité bayésienne a priori uniforme).
  \end{itemize}
\end{successbox}

% ============================================================
%  CHAPITRE 5 — RÉALISATION DE L'INTERFACE
% ============================================================
\chapter{Réalisation de l'Interface Web}
\label{chap:interface}

\section{Architecture de l'Application Django}

L'application web suit l'architecture MTV (Model–Template–View) de
Django~:

\begin{table}[h]
  \centering
  \caption{Structure du projet Django}
  \label{tab:structure}
  \begin{tabularx}{\textwidth}{llX}
    \toprule
    \textbf{Fichier} & \textbf{Rôle} & \textbf{Contenu clé} \\
    \midrule
    \fichier{views.py}           & Contrôleur & Prédiction, prétraitement, contexte \\
    \fichier{result.html}        & Vue         & Affichage verdict + jauges \\
    \fichier{settings.py}        & Config      & \code{MODEL\_PATH}, \code{STATS\_PATH} \\
    \fichier{brats\_classifier.py} & Entraînement & ResNet50, Loss, Youden \\
    \fichier{dataset\_stats.json} & Persistance & mean, std, threshold \\
    \bottomrule
  \end{tabularx}
\end{table}

% -----------------------------------------------------------
%  FIGURE 1 — Interface d'accueil / Drag & Drop
% -----------------------------------------------------------
\subsection{Interface d'Accueil}

La page d'accueil de BrainScan AI propose une zone de dépôt
(\textit{Drag \& Drop}) permettant à l'utilisateur de soumettre
directement son fichier IRM au format NIfTI (\fichier{.nii} ou
\fichier{.nii.gz}) sans installation locale d'aucun logiciel.

\begin{figure}[H]
  \centering
  % --> Remplacez 'fig_accueil' par le chemin réel de votre capture
  %     (ex. : screenshots/fig_accueil.png)
  \includegraphics[width=0.88\linewidth]{fig_accueil}
  \caption{%
    Interface d'accueil de \textbf{BrainScan AI} — Zone de dépôt
    (\textit{Drag \& Drop}) pour le fichier NIfTI~(\code{.nii.gz}).
    La charte graphique sombre avec l'accent bleu~(\code{\#00b4d8})
    confère à l'application un aspect médical et professionnel.
    L'utilisateur n'a aucun logiciel à installer~: tout le traitement
    s'effectue côté serveur Django.
  }
  \label{fig:accueil}
\end{figure}

\section{Flux de Prédiction Complet}

\begin{lstlisting}[caption={Vue de prédiction — flux principal (views.py)},
                   label={lst:predict_flow}]
def predict(request):
    # 1. Reception et sauvegarde du fichier .nii.gz
    tmp_path = sauvegarder_fichier(request.FILES["mri_file"])

    # 2. Extraction de la coupe axiale centrale
    slice_2d = _extract_central_slice(tmp_path)

    # 3. Pretraitement : Z-Score + CLAHE + resize + ImageNet
    tensor = _preprocess(slice_2d)

    # 4. Inference + inversion des classes
    raw = _get_model().predict(tensor, verbose=0)
    tumor_prob = 1.0 - float(raw[0][0])   # inversion 0=Tumeur

    # 5. Seuil optimal depuis dataset_stats.json
    stats = _get_stats()
    DECISION_THRESHOLD = float(stats.get("threshold", 0.50))
    has_tumor = tumor_prob >= DECISION_THRESHOLD

    # 6. Detection zone de signal atypique (80-95 %)
    zone_signal_atypique = (not has_tumor) and (0.80 <= tumor_prob <= 0.95)

    # 7. Rendu du template avec contexte enrichi
    return render(request, "detector/result.html", context)
\end{lstlisting}

% -----------------------------------------------------------
%  FIGURE 2 — Chargement et extraction de la coupe axiale
% -----------------------------------------------------------

Une fois le fichier soumis, le serveur effectue en séquence~:
extraction de la coupe axiale n°75, normalisation Z-Score
dataset-level, CLAHE, resize~$224\times224$, puis passage dans
le réseau ResNet50. La Figure~\ref{fig:chargement} illustre
l'état intermédiaire de l'interface pendant ce traitement.

\begin{figure}[H]
  \centering
  % --> Remplacez 'fig_chargement' par le chemin réel de votre capture
  \includegraphics[width=0.88\linewidth]{fig_chargement}
  \caption{%
    \textbf{Processus de chargement et d'extraction} — Après soumission
    du fichier \fichier{.nii.gz}, l'application extrait automatiquement
    la coupe axiale centrale (coupe~n°75, ou la coupe médiane si le
    volume compte moins de~75 coupes). La coupe brute est sauvegardée
    en PNG pour l'affichage côté client, tandis que la version
    prétraitée (Z-Score + CLAHE + normalisation ImageNet) est
    transmise au modèle ResNet50 pour l'inférence.
  }
  \label{fig:chargement}
\end{figure}

\section{Conception de l'Expérience Utilisateur}
\label{sec:ux}

\subsection{Philosophie de Design}

Le principal défi UX d'un outil de diagnostic assisté est d'informer
sans alarmer. Trois principes ont guidé les choix d'interface~:

\begin{enumerate}
  \item \textbf{Masquer les chiffres bruts inquiétants}~: une probabilité
        de «~79.3~\% de tumeur~» sur un verdict \textit{sain} crée une
        anxiété injustifiée. L'utilisateur final n'est pas statisticien.

  \item \textbf{Substituer la probabilité par la fiabilité}~: l'indicateur
        affiché répond à la question «~le modèle est-il sûr de son
        verdict~?~», pas «~quelle est la probabilité exacte~?~».

  \item \textbf{Contextualiser les cas incertains}~: si le signal tumoral
        est entre 80~\% et 95~\% sans franchir le seuil, un message
        spécifique oriente vers une consultation complémentaire.
\end{enumerate}

\subsection{Les Trois États d'Affichage}

\begin{table}[h]
  \centering
  \caption{États d'affichage selon la probabilité tumorale}
  \label{tab:etats_affichage}
  \begin{tabularx}{\textwidth}{lcXl}
    \toprule
    \textbf{État} & \textbf{Condition} & \textbf{Message} & \textbf{Couleur} \\
    \midrule
    Tumeur détectée
      & $p \geq \tau^*$
      & Probabilité tumorale brute
      & Rouge \\
    Signal atypique
      & $0.80 \leq p < \tau^*$ si $\tau^* > 0.80$
      & «~Analyse incertaine~»
      & Orange \\
    Cerveau sain
      & $p < 0.80$
      & «~Fiabilité~: Élevée — 100~\%~»
      & Vert \\
    \bottomrule
  \end{tabularx}
\end{table}

% -----------------------------------------------------------
%  FIGURES 3 & 4 — Résultats Sain et Tumeur (côte à côte)
% -----------------------------------------------------------

Les Figures~\ref{fig:resultat_sain} et~\ref{fig:resultat_tumeur}
illustrent les deux verdicts possibles renvoyés par BrainScan AI.
Chaque page de résultats présente~: la coupe IRM extraite (colonne
gauche), la jauge demi-cercle SVG animée, et l'indicateur
de fiabilité ou d'alerte (colonne droite).

\begin{figure}[H]
  \centering

  % ── Sous-figure gauche : Cerveau sain ────────────────────
  \begin{subfigure}[t]{0.48\textwidth}
    \centering
    % --> Remplacez 'fig_resultat_sain' par votre capture
    \includegraphics[width=\linewidth]{fig_resultat_sain}
    \caption{%
      \textbf{Verdict : Cerveau sain.}
      L'indicateur de \textit{Fiabilité du diagnostic} affiche
      \textbf{100\,\%} en vert~(\code{\#2ecc9a}). Le message
      rassurant \og Le modèle n'a détecté aucun signal dépassant
      le seuil de rigueur médicale \fg{} remplace l'affichage
      de la probabilité brute, évitant toute anxiété injustifiée.
    }
    \label{fig:resultat_sain}
  \end{subfigure}
  \hfill
  % ── Sous-figure droite : Tumeur détectée ─────────────────
  \begin{subfigure}[t]{0.48\textwidth}
    \centering
    % --> Remplacez 'fig_resultat_tumeur' par votre capture
    \includegraphics[width=\linewidth]{fig_resultat_tumeur}
    \caption{%
      \textbf{Verdict : Tumeur détectée.}
      L'alerte rouge~(\code{\#ff6b7a}) est déclenchée dès que
      $p_{\text{tumeur}} \geq \tau^*$. La jauge demi-cercle
      affiche le score de confiance du modèle. La coupe IRM
      est présentée à gauche pour permettre une vérification
      visuelle immédiate par le clinicien.
    }
    \label{fig:resultat_tumeur}
  \end{subfigure}

  \caption{%
    \textbf{Pages de résultats de BrainScan AI.}
    À gauche~: verdict \textit{Cerveau sain} avec indicateur de
    fiabilité vert à 100\,\%. À droite~: verdict \textit{Tumeur
    détectée} avec alerte rouge et score de confiance. Dans les
    deux cas, la coupe IRM axiale extraite est affichée pour
    permettre au clinicien une vérification visuelle du signal
    analysé.
  }
  \label{fig:resultats}
\end{figure}

\subsection{Implémentation Django Template}

\begin{lstlisting}[caption={Affichage conditionnel dans result.html (Django)},
                   language=HTML,
                   label={lst:template_ux}]
{% if zone_signal_atypique %}
  <!-- Zone orange : signal entre 80-95 % -->
  <div class="confidence-label">
    <span>Indice de confiance</span>
    <span style="color:#f39c12; font-weight:700;">Moderee</span>
  </div>
  <p style="color:#c8a84b; background:rgba(243,156,18,0.08);">
    <strong>Analyse incertaine : signaux atypiques.</strong>
    Signal detecte a {{ prob_tumeur }}% -- examen complementaire recommande.
  </p>
{% else %}
  <!-- Zone verte : haute fiabilite -->
  <div class="confidence-label">
    <span>Fiabilite du diagnostic</span>
    <span style="color:#2ecc9a; font-weight:700;">Elevee</span>
  </div>
  <p>Le modele n'a detecte aucun signal depassant le seuil
  de rigueur medicale.</p>
{% endif %}
\end{lstlisting}

\subsection{Jauge SVG Animée}

La jauge demi-cercle est rendue en SVG pur, animée par JavaScript
avec une transition CSS de 1.2 secondes~:

\begin{lstlisting}[caption={Animation de la jauge SVG (result.html — JS)},
                   language=JavaScript,
                   label={lst:gauge_js}]
const hasTumor          = {{ has_tumor|yesno:"true,false" }};
const zoneSignalAtypique = {{ zone_signal_atypique|yesno:"true,false" }};

window.addEventListener('load', () => {
  const arc   = document.getElementById('gauge-arc');
  const total = 251.2;   // pi * 80 (rayon SVG)

  const pct    = (hasTumor || zoneSignalAtypique) ? tumorPct : 100;
  const filled = (pct / 100) * total;

  setTimeout(() => {
    arc.setAttribute('stroke-dasharray',
      filled + ' ' + (total - filled));
  }, 300);
});
\end{lstlisting}

\section{Chargement Paresseux du Modèle}

Pour éviter de recharger le modèle ResNet50 à chaque requête
(coût~: 1–3 secondes), un mécanisme de singleton est utilisé~:

\begin{lstlisting}[caption={Singleton de chargement du modèle (views.py)},
                   label={lst:lazy_load}]
_model = None

def _get_model():
    global _model
    if _model is None:
        focal_smooth_loss = _build_focal_smooth_loss(
            gamma=2.0, alpha=0.25, label_smoothing=0.1
        )
        _model = load_model(
            str(settings.MODEL_PATH),
            custom_objects={"focal_smooth_loss": focal_smooth_loss},
        )
    return _model
\end{lstlisting}

% ============================================================
%  CHAPITRE 6 — CONCLUSION ET PERSPECTIVES
% ============================================================
\chapter{Conclusion et Perspectives}
\label{chap:conclusion}

\section{Bilan}

Ce projet a abouti à un système complet de détection de tumeurs
cérébrales, couvrant l'ensemble de la chaîne~: collecte et prétraitement
des données IRM, entraînement d'un modèle ResNet50 avec une fonction de
perte adaptée, calibration automatique du seuil de décision, et
déploiement via une interface web Django soignée.

Les principaux apports techniques sont~:

\begin{enumerate}
  \item Un pipeline de normalisation Z-Score \textit{dataset-level}
        garantissant la cohérence entre entraînement et inférence.
  \item Une fonction de perte \textit{Focal Loss + Label Smoothing}
        réduisant la surconfiance du modèle sur les données médicales.
  \item Un calcul automatique du seuil de Youden évitant toute
        configuration manuelle post-déploiement.
  \item Une correction robuste de l'inversion des classes par simple
        complémentation de la probabilité brute.
  \item Une interface UX orientée \textit{réassurance} masquant les
        probabilités brutes anxiogènes.
\end{enumerate}

\section{Perspectives}

\subsection{Améliorations Techniques}

\begin{itemize}
  \item \textbf{Segmentation fine}~: migrer vers une architecture
        U-Net pour produire un masque de localisation de la tumeur,
        en plus du verdict binaire.

  \item \textbf{Multi-modalité}~: exploiter les quatre séquences IRM
        disponibles dans BraTS (T1, T1ce, T2, FLAIR) au lieu d'une
        seule coupe FLAIR, par fusion de features.

  \item \textbf{Calibration de probabilité}~: appliquer une régression
        isotonique ou un Temperature Scaling post-entraînement pour
        que \(p_{\text{tumeur}} = 0.7\) signifie réellement
        «~70~\% de cas tumoraux dans la population~».

  \item \textbf{Grad-CAM}~: intégrer la visualisation par carte de
        saillance pour localiser visuellement les zones d'attention
        du modèle sur l'IRM.
\end{itemize}

\subsection{Vers un Déploiement Clinique}

Un déploiement en milieu clinique nécessiterait~:
\begin{itemize}
  \item Une validation prospective sur une cohorte externe (hôpital
        partenaire, données non vues).
  \item Un audit de biais~: vérifier que le modèle est équitable sur
        tous les groupes démographiques (âge, sexe, ethnicit").
  \item Une certification médicale (marquage CE classe IIb en Europe,
        clearance FDA 510(k) aux États-Unis).
\end{itemize}

\vspace{1cm}

\begin{successbox}[Mot de fin]
  Ce projet illustre comment les techniques de \textit{deep learning}
  modernes peuvent être mises au service d'un enjeu de santé publique,
  à condition d'être accompagnées d'une réflexion rigoureuse sur la
  calibration, la robustesse et l'expérience utilisateur.
\end{successbox}

% ============================================================
%  BIBLIOGRAPHIE
% ============================================================
\printbibliography[heading=bibintoc, title={Références Bibliographiques}]

% ============================================================
%  ANNEXES
% ============================================================
\appendix

\chapter{Configuration du Projet Django}
\label{ann:settings}

\begin{lstlisting}[caption={Extraits de settings.py},
                   label={lst:settings}]
from pathlib import Path

BASE_DIR = Path(__file__).resolve().parent.parent

MODEL_PATH = BASE_DIR / "brats_final.keras"
STATS_PATH = BASE_DIR / "dataset_stats.json"

DATA_UPLOAD_MAX_MEMORY_SIZE = 52_428_800   # 50 Mo
FILE_UPLOAD_MAX_MEMORY_SIZE = 52_428_800
\end{lstlisting}

\chapter{Format du Fichier dataset\_stats.json}
\label{ann:json}

\begin{lstlisting}[caption={dataset\_stats.json complet après entraînement},
                   label={lst:json_full}]
{
  "mean":      1729.701416015625,
  "std":       16289.5810546875,
  "n_images":  875,
  "threshold": 0.5312
}
\end{lstlisting}

La clé \code{"threshold"} est absente si \code{find\_optimal\_threshold()}
n'a pas encore été appelée. Dans ce cas, \fichier{views.py} utilise
le fallback \(0.50\).

\chapter{Hyperparamètres d'Entraînement}
\label{ann:hyperparams}

\begin{table}[h]
  \centering
  \caption{Hyperparamètres du modèle final}
  \label{tab:hyperparams}
  \begin{tabular}{ll}
    \toprule
    \textbf{Paramètre} & \textbf{Valeur} \\
    \midrule
    Architecture           & ResNet50 (ImageNet) \\
    Couches fine-tunées    & 20 dernières \\
    Taille d'entrée        & $224 \times 224 \times 3$ \\
    Taille de batch        & 16 \\
    Époques maximales      & 50 \\
    Taux d'apprentissage   & $10^{-4}$ (Adam) \\
    L2 régularisation      & $10^{-3}$ \\
    Dropout                & $0.50$ (deux couches) \\
    Label smoothing        & $\varepsilon = 0.10$ \\
    Focal Loss $\gamma$    & $2.0$ \\
    Focal Loss $\alpha$    & $0.25$ \\
    Poids classe tumeur    & $2.0$ \\
    EarlyStopping          & patience = 5 (val\_loss) \\
    ReduceLROnPlateau      & factor = 0.5, patience = 3 \\
    Seuil de décision      & Youden, clamped $[0.25, 0.75]$ \\
    \bottomrule
  \end{tabular}
\end{table}

\end{document}
